\chapter{Обзор литературы}
\label{chap:literature}
\setlength{\parskip}{0.8em}

\section{Клинический обзор ожоговой болезни}

Ожоговые травмы являются одной из наиболее распространённых причин травматизма во всём мире. По данным Всемирной организации здравоохранения, ежегодно от ожогов получают травмы около 11 миллионов человек, при этом около 180~000 случаев заканчиваются летальным исходом~\cite{who2018burns}. Большинство смертей происходит в странах с низким и средним уровнем дохода, где доступ к специализированной медицинской помощи ограничен.

Ожоги классифицируются по глубине поражения тканей на три основные степени~\cite{peck2012burns}:

\textbf{Ожоги первой степени (поверхностные)} затрагивают только эпидермис — верхний слой кожи. Характеризуются покраснением, болью и лёгким отёком. Заживают самостоятельно в течение 3–7 дней без образования рубцов. Примеры: солнечные ожоги, кратковременный контакт с горячей поверхностью.

\textbf{Ожоги второй степени (частичной толщины)} проникают в дерму. Разделяются на поверхностные и глубокие. Поверхностные затрагивают папиллярную дерму, характеризуются образованием волдырей, выраженной болью, влажной поверхностью и заживают за 2–3 недели. Глубокие проникают в ретикулярную дерму, могут требовать хирургического лечения.

\textbf{Ожоги третьей степени (полной толщины)} разрушают все слои кожи, включая нервные окончания, поэтому они могут быть безболезненны. Характеризуются белым, коричневым или чёрным цветом, сухой кожистой текстурой. Требуют хирургического вмешательства — некрэктомии и пересадки кожи.

\section{Методы оценки глубины ожогов}

\subsection{Клиническая оценка}

Традиционно глубина ожога оценивается визуально опытным специалистом на основании цвета, текстуры, чувствительности и наличия волдырей. Однако точность клинической оценки варьируется от 50\% до 75\%~\cite{monstrey2008assessment}, особенно в первые 48–72~часа после травмы, когда ожог может прогрессировать.

\subsection{Инструментальные методы}

Среди инструментальных методов применяются лазерная доплеровская визуализация (LDI), термография и оптическая когерентная томография~\cite{hoeksema2009accuracy}. Эти методы имеют высокую точность ($>$95\%), но требуют дорогостоящего оборудования и квалифицированного персонала, что ограничивает их применение.

\section{Компьютерное зрение в медицинской диагностике}

\subsection{Свёрточные нейронные сети}

Свёрточные нейронные сети (CNN), впервые описанные LeCun et al.~\cite{lecun1998cnn}, совершили революцию в анализе медицинских изображений. Архитектура CNN извлекает иерархические признаки из изображений через последовательность свёрточных слоёв:

\begin{equation}
    \mathbf{h}^{(l)} = \sigma\left(\mathbf{W}^{(l)} * \mathbf{h}^{(l-1)} + \mathbf{b}^{(l)}\right)
\label{eq:cnn}
\end{equation}

\noindent где $\mathbf{h}^{(l)}$ — выход слоя $l$; $\mathbf{W}^{(l)}$ — ядро свёртки; $\mathbf{b}^{(l)}$ — смещение; $\sigma(\cdot)$ — функция активации (обычно ReLU).

\subsection{Трансферное обучение}

Трансферное обучение~\cite{yosinski2014transfer} позволяет использовать предобученные на больших датасетах (ImageNet, COCO) модели для специализированных медицинских задач. Нижние слои сети, извлекающие общие признаки (края, текстуры), замораживаются, а верхние слои дообучаются на целевом датасете.

\subsection{Применение CNN в дерматологии}

Esteva et al.~\cite{esteva2017dermatologist} продемонстрировали, что CNN может классифицировать рак кожи на уровне дерматолога. Подобные работы применялись для классификации псориаза, экземы и других кожных заболеваний. В Таблице~\ref{tab:lit_review} приведён обзор избранных исследований.

\begin{table}[H]
\centering
\caption{Обзор исследований по применению глубокого обучения в анализе кожных поражений}
\label{tab:lit_review}
\begin{tabular}{p{3.5cm} p{3cm} p{2.5cm} p{2cm} p{2cm}}
\toprule
\textbf{Исследование} & \textbf{Задача} & \textbf{Модель} & \textbf{Accuracy} & \textbf{Датасет} \\
\midrule
Esteva et al.\ (2017)~\cite{esteva2017dermatologist} & Рак кожи & Inception v3 & 91\% & 130K изобр. \\
Chauhan et al.\ (2020)~\cite{chauhan2020burn} & Ожоги & ResNet-50 & 85\% & 1200 изобр. \\
Abubakar et al.\ (2020)~\cite{abubakar2020burn} & Ожоги & VGG-16 & 82\% & 500 изобр. \\
Khan et al.\ (2021)~\cite{khan2021burn} & Ожоги & DenseNet & 87\% & 2000 изобр. \\
\textbf{Данная работа} & \textbf{Детекция ожогов} & \textbf{YOLOv8n} & \textbf{mAP 0.85} & \textbf{1225 изобр.} \\
\bottomrule
\end{tabular}
\end{table}

\section{Эволюция архитектуры YOLO}

\subsection{YOLO: You Only Look Once}

Архитектура YOLO (You Only Look Once), представленная Redmon et al.~\cite{redmon2016yolo}, предложила принципиально новый подход к детекции объектов — однопроходную архитектуру, обрабатывающую изображение целиком и выдающую предсказания за один проход через сеть. В отличие от двухэтапных методов (R-CNN, Fast R-CNN), YOLO обеспечивает детекцию в реальном времени.

Основная идея: изображение делится на сетку $S \times S$, каждая ячейка предсказывает $B$ ограничивающих прямоугольников и вероятности классов:

\begin{equation}
    \text{Prediction} = (x, y, w, h, \text{confidence}, p_{c_1}, p_{c_2}, \ldots, p_{c_n})
\label{eq:yolo}
\end{equation}

\noindent где $(x, y)$ — центр бокса; $(w, h)$ — ширина и высота; confidence — уверенность в наличии объекта; $p_{c_i}$ — вероятность класса $i$.

\subsection{Эволюция версий}

\textbf{YOLOv2/YOLO9000}~\cite{redmon2017yolo9000} добавил batch normalization, якорные боксы и обучение на нескольких разрешениях.

\textbf{YOLOv3}~\cite{redmon2018yolov3} использовал Darknet-53 backbone и детекцию на трёх масштабах через Feature Pyramid Network (FPN).

\textbf{YOLOv4}~\cite{bochkovskiy2020yolov4} внёс CSPDarknet backbone, PANet, Mish activation и множество техник аугментации (Mosaic, CutMix).

\textbf{YOLOv5} (Ultralytics) — PyTorch-реализация с автоматическими якорями и поддержкой экспорта в различные форматы.

\textbf{YOLOv8}~\cite{ultralytics2023yolov8} — последняя версия от Ultralytics с anchor-free архитектурой, C2f-модулями и унифицированным API для detection, segmentation и pose estimation.

\section{Архитектура YOLOv8}

YOLOv8 представляет собой современную anchor-free архитектуру детекции объектов. Ключевые компоненты:

\textbf{Backbone} (CSPDarknet) извлекает признаки из входного изображения через последовательность свёрточных блоков C2f (Cross Stage Partial с двумя convolutions), использующих остаточные соединения.

\textbf{Neck} (PANet — Path Aggregation Network) агрегирует признаки разных масштабов, обеспечивая детекцию объектов разного размера.

\textbf{Head} (Decoupled Head) разделяет классификацию и регрессию бокса на независимые ветви, что улучшает качество предсказаний.

YOLOv8 доступен в пяти вариантах по размеру: n (nano), s (small), m (medium), l (large), x (extra-large). YOLOv8n, использованный в данной работе, имеет около 3,2~млн параметров и обеспечивает баланс между скоростью и точностью.

Функция потерь YOLOv8 состоит из трёх компонентов:

\begin{equation}
    \mathcal{L} = \lambda_{box} \mathcal{L}_{box} + \lambda_{cls} \mathcal{L}_{cls} + \lambda_{dfl} \mathcal{L}_{dfl}
\label{eq:yolov8loss}
\end{equation}

\noindent где $\mathcal{L}_{box}$ — потери регрессии бокса (CIoU); $\mathcal{L}_{cls}$ — потери классификации (Binary Cross-Entropy); $\mathcal{L}_{dfl}$ — Distribution Focal Loss для уточнения границ.

\section{Метрики оценки детекции объектов}

\subsection{Precision и Recall}

Precision (точность) — доля правильных детекций среди всех детекций:

\begin{equation}
    P = \frac{TP}{TP + FP}
\label{eq:precision}
\end{equation}

Recall (полнота) — доля обнаруженных объектов среди всех существующих:

\begin{equation}
    R = \frac{TP}{TP + FN}
\label{eq:recall}
\end{equation}

\subsection{IoU (Intersection over Union)}

IoU измеряет степень перекрытия предсказанного бокса с ground truth:

\begin{equation}
    \text{IoU} = \frac{|B_{pred} \cap B_{gt}|}{|B_{pred} \cup B_{gt}|}
\label{eq:iou}
\end{equation}

Детекция считается true positive, если $\text{IoU} \geq \text{threshold}$ (обычно 0,5 или 0,5:0,95).

\subsection{mAP (mean Average Precision)}

Average Precision вычисляется как площадь под кривой precision-recall. mAP — среднее AP по всем классам:

\begin{equation}
    \text{mAP} = \frac{1}{|\mathcal{C}|} \sum_{c \in \mathcal{C}} \text{AP}_c
\label{eq:map_def}
\end{equation}

mAP@0.5 использует порог IoU = 0,5; mAP@0.5:0.95 усредняет по порогам от 0,5 до 0,95 с шагом 0,05.

\section{Выводы и пробелы в исследованиях}

Литература устанавливает:
\begin{enumerate}
    \item Глубокое обучение, особенно CNN, демонстрирует высокую эффективность в анализе медицинских изображений, включая дерматологические задачи.
    \item Архитектура YOLO обеспечивает детекцию объектов в реальном времени с конкурентной точностью.
    \item Существующие работы по классификации ожогов преимущественно используют классификационные CNN (ResNet, VGG), не предоставляя локализацию области поражения.
    \item Применение YOLOv8 для задачи детекции и классификации ожоговых поражений остаётся малоизученной областью.
\end{enumerate}

Данная диссертация устраняет эти пробелы путём применения современной архитектуры YOLOv8 для детекции и классификации ожогов с одновременной локализацией области поражения.
